% GENERATED FILE. Edit canonical lessons, then run npm run generate:books.
% canonical-source-hash: fnv1a64:bd44eb10e3ae6e73
% canonical-lessons: AR-C532-walid, AR-C532-rabib, AR-C532-qaraba, AR-C532-asha, AR-C532-mata

\chapter{A life and its end}
\label{ch:ar-a-life-and-its-end}
\hlchaptermodality{532}

\begin{chapteropening}
\textbf{By the end of this chapter:} \emph{I can talk about a parent, a stepson and kinship, and say that someone lived a life or died.}

Complete the last lesson of chapter 532: I can talk about a parent, a stepson and kinship, and say that someone lived a life or died.
\end{chapteropening}

\section[\texorpdfstring{\ar{والد} (wālid)}{والد (wālid)}]{\ar{والد} (wālid) — a parent; a father}

\label{lesson:AR-C532-walid}



\begin{quote}
Before the new one: say the Arabic for old age, then the Arabic for a nanny.
\end{quote}

\subsection*{You'll want to know: \ar{والد}}
\textbf{\ar{والد}} — \emph{wālid} — \textquotedblleft{}a parent; a father\textquotedblright{}.

\begin{quote}
\ar{سافر} \ar{والدي} \ar{إلى} \ar{المدينة}.
\end{quote}

\begin{quote}
\emph{sāfara wālidī ilā al-madīna.}
\end{quote}

\begin{quote}
My father travelled to the city.
\end{quote}

\ar{الوالدان} means the parents; \ar{والدة} is a mother.

\begin{cousinweb}
\textbf{\ar{والد}} (\emph{wālid}) is built on the native Arabic root \emph{w-l-d}, \textquotedblleft{}to give birth\textquotedblright{} (Lane). Walad, a boy, and mīlād, birth, which you know, share the root.
\end{cousinweb}

\begin{grammarlens}[title={What you've built}]
One more word for describing what you lived and felt.
\end{grammarlens}

\subsection*{Your turn}
\begin{itemize}
\raggedright
  \item \emph{Say it:} \emph{wālid}
  \item \emph{Say it:} \emph{wālid}, once more
  \item \emph{Say it:} \emph{sāfara wālidī ilā al-madīna.} — My father travelled to the city.
  \item \emph{Recall:} say the Arabic for old age, then the Arabic for a nanny, then say \emph{wālid} again
\end{itemize}

\begin{tcolorbox}[breakable,colback=teal!4,colframe=teal!35!black,title={Before you move on}]
What does \ar{والد} mean? (\textquotedblleft{}A parent; a father\textquotedblright{}.) Say it once more.
\end{tcolorbox}
\section[\texorpdfstring{\ar{ربيب} (rabīb)}{ربيب (rabīb)}]{\ar{ربيب} (rabīb) — a stepson}

\label{lesson:AR-C532-rabib}



\begin{quote}
Before the new one: say the Arabic for a nanny, then the Arabic for a parent.
\end{quote}

\subsection*{You'll want to know: \ar{ربيب}}
\textbf{\ar{ربيب}} — \emph{rabīb} — \textquotedblleft{}a stepson\textquotedblright{}.

\begin{quote}
\ar{يحب} \ar{الرجل} \ar{ربيبه} \ar{مثل} \ar{ابنه}.
\end{quote}

\begin{quote}
\emph{yuḥibbu ar-rajul rabībahu mithla ibnihi.}
\end{quote}

\begin{quote}
The man loves his stepson like his own son.
\end{quote}

\ar{ربيبة} is a stepdaughter.

\begin{cousinweb}
\textbf{\ar{ربيب}} (\emph{rabīb}) is built on the native Arabic root \emph{r-b-b}, \textquotedblleft{}a lord, an owner; to rear\textquotedblright{} (Lane). Lane gives rabīb as a wife's son, reared by her husband.
\end{cousinweb}

\begin{grammarlens}[title={What you've built}]
One more word for describing what you lived and felt.
\end{grammarlens}

\subsection*{Your turn}
\begin{itemize}
\raggedright
  \item \emph{Say it:} \emph{rabīb}
  \item \emph{Say it:} \emph{rabīb}, once more
  \item \emph{Say it:} \emph{yuḥibbu ar-rajul rabībahu mithla ibnihi.} — The man loves his stepson like his own son.
  \item \emph{Recall:} say the Arabic for a nanny, then the Arabic for a parent, then say \emph{rabīb} again
\end{itemize}

\begin{tcolorbox}[breakable,colback=teal!4,colframe=teal!35!black,title={Before you move on}]
What does \ar{ربيب} mean? (\textquotedblleft{}A stepson\textquotedblright{}.) Say it once more.
\end{tcolorbox}
\section[\texorpdfstring{\ar{قرابة} (qarāba)}{قرابة (qarāba)}]{\ar{قرابة} (qarāba) — kinship}

\label{lesson:AR-C532-qaraba}



\begin{quote}
Before the new one: say the Arabic for a parent, then the Arabic for a stepson.
\end{quote}

\subsection*{You'll want to know: \ar{قرابة}}
\textbf{\ar{قرابة}} — \emph{qarāba} — \textquotedblleft{}kinship\textquotedblright{}.

\begin{quote}
\ar{بيني} \ar{وبين} \ar{جاري} \ar{قرابة} \ar{بعيدة}.
\end{quote}

\begin{quote}
\emph{baynī wa-bayna jārī qarāba baʿīda.}
\end{quote}

\begin{quote}
My neighbour and I are distant relatives.
\end{quote}

\begin{cousinweb}
\textbf{\ar{قرابة}} (\emph{qarāba}) is built on the native Arabic root \emph{q-r-b}, \textquotedblleft{}to be near\textquotedblright{} (Lane). Qarīb, near and a relative, which you know, is from the same root.
\end{cousinweb}

\begin{grammarlens}[title={What you've built}]
One more word for describing what you lived and felt.
\end{grammarlens}

\subsection*{Your turn}
\begin{itemize}
\raggedright
  \item \emph{Say it:} \emph{qarāba}
  \item \emph{Say it:} \emph{qarāba}, once more
  \item \emph{Say it:} \emph{baynī wa-bayna jārī qarāba baʿīda.} — My neighbour and I are distant relatives.
  \item \emph{Recall:} say the Arabic for a parent, then the Arabic for a stepson, then say \emph{qarāba} again
\end{itemize}

\begin{tcolorbox}[breakable,colback=teal!4,colframe=teal!35!black,title={Before you move on}]
What does \ar{قرابة} mean? (\textquotedblleft{}Kinship\textquotedblright{}.) Say it once more.
\end{tcolorbox}
\section[\texorpdfstring{\ar{عاش} (ʿāsha)}{عاش (ʿāsha)}]{\ar{عاش} (ʿāsha) — to lead a life}

\label{lesson:AR-C532-asha}



\begin{quote}
Before the new one: say the Arabic for a stepson, then the Arabic for kinship.
\end{quote}

\subsection*{You'll want to know: \ar{عاش}}
\textbf{\ar{عاش}} — \emph{ʿāsha} — \textquotedblleft{}to lead a life\textquotedblright{}.

\begin{quote}
\ar{عاش} \ar{جدي} \ar{حياة} \ar{بسيطة} \ar{في} \ar{الريف}.
\end{quote}

\begin{quote}
\emph{ʿāsha jaddī ḥayāh basīṭa fī ar-rīf.}
\end{quote}

\begin{quote}
My grandfather lived a simple life in the countryside.
\end{quote}

Present: \ar{يعيش}, \emph{yaʿīsh}.

\begin{cousinweb}
\textbf{\ar{عاش}} (\emph{ʿāsha}) is built on the native Arabic root \emph{ʿ-y-sh}, \textquotedblleft{}to live\textquotedblright{} (Lane). The noun ʿaysh is life and, in Egypt, bread.
\end{cousinweb}

\begin{grammarlens}[title={What you've built}]
One more word for describing what you lived and felt.
\end{grammarlens}

\subsection*{Your turn}
\begin{itemize}
\raggedright
  \item \emph{Say it:} \emph{ʿāsha}
  \item \emph{Say it:} \emph{ʿāsha}, once more
  \item \emph{Say it:} \emph{ʿāsha jaddī ḥayāh basīṭa fī ar-rīf.} — My grandfather lived a simple life in the countryside.
  \item \emph{Recall:} say the Arabic for a stepson, then the Arabic for kinship, then say \emph{ʿāsha} again
\end{itemize}

\begin{tcolorbox}[breakable,colback=teal!4,colframe=teal!35!black,title={Before you move on}]
What does \ar{عاش} mean? (\textquotedblleft{}To lead a life\textquotedblright{}.) Say it once more.
\end{tcolorbox}
\section[\texorpdfstring{\ar{مات} (māta)}{مات (māta)}]{\ar{مات} (māta) — to die}

\label{lesson:AR-C532-mata}



\begin{quote}
Before the new one: say the Arabic for kinship, then the Arabic for to lead a life.
\end{quote}

\subsection*{You'll want to know: \ar{مات}}
\textbf{\ar{مات}} — \emph{māta} — \textquotedblleft{}to die\textquotedblright{}.

\begin{quote}
\ar{مات} \ar{الكلب} \ar{العجوز} \ar{في} \ar{الشتاء}.
\end{quote}

\begin{quote}
\emph{māta al-kalb al-ʿajūz fī ash-shitāʾ.}
\end{quote}

\begin{quote}
The old dog died in the winter.
\end{quote}

Present: \ar{يموت}, \emph{yamūt}.

\ar{توفي} is the polite word for a person.

\begin{cousinweb}
\textbf{\ar{مات}} (\emph{māta}) is built on the native Arabic root \emph{m-w-t}, \textquotedblleft{}to die\textquotedblright{} (Lane). Mayyit, dead, which you know, is from the same root.
\end{cousinweb}

\begin{grammarlens}[title={What you've built}]
One more word for describing what you lived and felt.
\end{grammarlens}

\subsection*{Your turn}
\begin{itemize}
\raggedright
  \item \emph{Say it:} \emph{māta}
  \item \emph{Say it:} \emph{māta}, once more
  \item \emph{Say it:} \emph{māta al-kalb al-ʿajūz fī ash-shitāʾ.} — The old dog died in the winter.
  \item \emph{Recall:} say the Arabic for kinship, then the Arabic for to lead a life, then say \emph{māta} again
\end{itemize}

\begin{tcolorbox}[breakable,colback=teal!4,colframe=teal!35!black,title={Before you move on}]
What does \ar{مات} mean? (\textquotedblleft{}To die\textquotedblright{}.) Say it once more.
\end{tcolorbox}
